\begin{table}[tbp]
\caption{What the archive study adds to the evaluation framework.}
\label{tab:ch8-feedback}
\centering\small\renewcommand{\arraystretch}{1.14}
\begin{tabularx}{\linewidth}{@{}>{\raggedright\arraybackslash}p{25mm}>{\raggedright\arraybackslash}X>{\raggedright\arraybackslash}X@{}}\toprule
Concept & Empirical anchor & Refinement and boundary\\\midrule
Evaluated object & Reasons include successors, relocation, restructuring, and replacement. & Follow the relevant capability across repositories; a destination link alone does not demonstrate continuity.\\
Observation & Reasons, practices, and trajectories use different subsamples. & Visibility is question-specific; missing explanations remain unknown.\\
Time and conditions & Long post-peak intervals and explicit transition schedules. & Separate lifecycle landmarks and horizons; a last commit is not an archival date.\\
Models and use & Commit concentration appears at the tree's first split. & Use the finding to examine maintenance responsibility; a split rule is not a universal risk threshold.\\
Value and consequences & Alternatives, forks, guides, and handover attempts. & Assess both retained capability and transferred burden; adoption of a practice does not establish its effects.\\\bottomrule
\end{tabularx}
\tablenote{Synthesis based on \citet[Chapter 4]{xia2024thesis} and the analyses in this chapter. It specifies how the evidence qualifies and develops Chapters 2--3.}
\end{table}
